\section{Introduction}
\label{sec:introduction}

Begin with the complete graph $K_n$. At each step, choose uniformly from
the triangles whose three edges are still present, and delete those edges.
Stop when no triangle remains. The removed triangles form an edge-disjoint
packing; the edges left over form its \emph{leave}. Let $F_n$ be the number
of edges in this terminal graph. The problem is to determine the typical
size of this leave, including its leading constant.

The accepted triangles form a partial Steiner triple system, and $F_n$
measures its uncovered pairs. Bollob\'as and Erd\H{o}s conjectured in 1990
that the expected leave has order $n^{3/2}$, as reported in
\cite[Introduction]{BohmanFriezeLubetzky2015}.
Spencer's branching-process analysis and the independent nibble argument
of R\"odl and Thoma showed that the leave is $o(n^2)$ with high probability
\cite{Spencer1995,RodlThoma1996}. Grable proved bounds
$O(n^{11/6+\xi})$ for each fixed $\xi>0$ and outlined the higher-order
calculations leading to $O(n^{7/4+\xi})$
\cite[Theorems~5 and~7]{Grable1997}.
Bohman, Frieze and Lubetzky gave a short proof of the bound
$O(n^{7/4}\log^{5/4}n)$ \cite[Theorem~1]{BohmanFriezeLubetzky2010}
and subsequently proved $F_n=n^{3/2+o(1)}$ with high probability
\cite[Theorem~1]{BohmanFriezeLubetzky2015}. Their analysis uses
self-correcting dynamic concentration and carefully chosen triangular
extensions to determine the exponent.

Joos and K\"uhn developed removal estimates for general strictly balanced
hypergraphs, including pseudorandom starting hypergraphs
\cite[Theorems~1.2--1.3]{joosKuehn2025}. Their Conjecture~16.2 predicts
an asymptotic constant for the terminal leave; its triangle specialization
is $1/(2\sqrt2)$ in probability. The following theorem proves this
triangle case with $L^2$ convergence. Numbered references to Joos and
K\"uhn refer throughout to version~2.

\begin{theorem}\label{thm:main}
For the uniform triangle-removal process starting from $K_n$,
\[
 \E\left[\left(\frac{F_n}{n^{3/2}}-
                   \frac1{2\sqrt2}\right)^2\right]\longrightarrow0.
\]
In particular,
\[
 \frac{F_n}{n^{3/2}}\xrightarrow{\Prb}\frac1{2\sqrt2},
 \qquad
 \frac{\E F_n}{n^{3/2}}\longrightarrow\frac1{2\sqrt2}.
\]
\end{theorem}

The only external proof input is the early-prefix control from
Joos and K\"uhn, specialized in Section~\ref{sec:prefix}. All subsequent
continuation estimates are proved here. The theorem concerns removal
from $K_n$; it does not assert a sharp constant for arbitrary starting
graphs or for general hypergraph removal, or a fluctuation law.

\subsection{The continuation argument}

The proof separates the part of the process governed by uniform subgraph
counts from the final evolution of individual edges. Stop at a deterministic
time when the edge density is approximately $p=n^{-1/2+\epsilon}$, for a
small fixed $\epsilon>0$. The graph then has
\[
 m=\frac{n^2p}{2}\quad\text{edges},\qquad D=np^2
\]
triangles through each edge to relative error $o(1)$.
Section~\ref{sec:prefix} states the exact properties needed and verifies
their application from the public removal estimates.

The continuation analysis uses backward priority exploration, whose
hypergraph form appears in Spencer's work
\cite[Sections~2--4]{Spencer1995}. Grable formulated triangle packing
as a uniform priority scan and used adaptive exposure decision trees in
his concentration analysis
\cite[Sections~1 and~4]{Grable1997}. The quantitative issue here is to
follow a growing rescaled time interval and to approximate joint survival
on the scale of its vanishing probability. A fixed-time local
approximation alone does not give the precision required for a second
moment.

Give the triangles of this graph fresh independent uniform priorities and
scan them in increasing order, accepting a triangle when all its edges
remain. This represents the continuation of uniform removal. An edge
survives precisely when a recursive test returns true. A child test asks
whether the other \emph{two} edges of a candidate triangle survive before
its priority. We test the two edges with one jointly ordered list of
candidates, which is essential to the later visitation estimate.
Section~\ref{sec:queries} proves this rule and constructs
an auxiliary tree in which each occurrence of a candidate receives an
independent priority.

In this independent model, the survival probabilities satisfy exact
product equations. The scalar comparison is
\[
 q(t)=(1+2Dt)^{-1/2},\qquad 0\le t\le1.
\]
This is Spencer's scalar branching law after the time change $c=Dt$
\cite[Section~3, equation~(9)]{Spencer1995}. Bal and Bennett obtained
related matching estimates for fixed-degree random regular hypergraphs
and regular hypergraphs of growing girth
\cite[Corollaries~3--4]{BalBennett2023}; here $D$ grows with $n$,
and the terminal rescaled time is $D$.
To prove a uniform approximation by $q$, one needs stability of a matrix
acting on the graph's edges. We obtain it from cycle counts in vertex
links. Each normalized link matrix is close in Euclidean norm to its
averaging projection; the sum of these projections has a directly
controlled semigroup in the maximum norm. A finite perturbation expansion
then gives the stability needed by the product equations
(Section~\ref{sec:stability}).

The final issue is the difference between independent occurrences and
repeated triangle types in the finite graph. We compare the complete
candidate lists exposed by one or two root queries. A repetition forces
an extra graph edge in the union of at most two ancestral query paths
(Section~\ref{sec:coupling}). Small rooted-template bounds control the
number of such witnesses. Joint time ordering controls the probability
that their calls are visited. Their combination gives an absolute error
$o(1/D)$ (Section~\ref{sec:witnesses}). This precision matters because the
two-root survival probability is itself of order $1/D$.

The one- and two-root probabilities determine the first and second
moments of $F_n$. Their normalization is
\[
 mq(1)=\frac{n^2p}{2\sqrt{1+2np^2}}
       \sim\frac{n^{3/2}}{2\sqrt2}.
\]
Uniformity over the good prefix graphs, together with a sufficiently
small prefix failure probability, gives the unconditional conclusion
in Section~\ref{sec:moments}.

The continuation estimates developed here are the maximum-norm stability
of the unfolding and the smaller than $D^{-1}$ error for adaptive paired
queries. Both the numerical constant and the moment conclusion use
these estimates.

\subsection{Conventions}

All graphs are finite and simple. An embedding is injective and preserves
edges; it need not preserve nonedges. Thus an embedding of a fixed labeled
cycle counts its ordered vertex labels, without division by automorphisms.
The notation $(1\pm a)b$ means a number in $[(1-a)b,(1+a)b]$.
All unspecified constants below are independent of $n$ and of the realized
good prefix graph. They may depend on the fixed parameters introduced
next. A statement holds with superpolynomially high probability if its
failure probability is $n^{-\omega(1)}$, meaning smaller than $n^{-C}$
for every fixed $C>0$ when $n$ is sufficiently large.
